\documentclass{article}
\usepackage{amsmath,amssymb,amsthm,xcolor}
\newcommand{\verifierissue}[1]{%
  \par\noindent\textcolor{red}{\textbf{[Verifier issue:]}\ #1}\par\medskip}


\newtheorem{conjecture}{Conjecture}

\title{Stabilizing quantities $\rho_p$ for boundary control of PDE}
\author{Amaury Hayat}
\date{February 09, 2026}

\begin{document}
\maketitle

\section*{Context}

This is a question originally raised by Jean-Michel Coron in 2016.
When studying the stability of nonlinear PDE of the form
\begin{equation}
\partial_t u + A(u)\partial_x u + B(u) = 0, \quad \text{on } [0, L],
\end{equation}
where $u \in \mathbb{R}^n$ and with boundary conditions of the form
\begin{equation}
\begin{pmatrix} u_+(t, 0) \\ u_-(t, L) \end{pmatrix}
= G \begin{pmatrix} u_+(t, L) \\ u_-(t, 0) \end{pmatrix}
\end{equation}
a natural quantity that appears is the following function
\begin{equation}
\rho_p : M \mapsto \inf\{\|\Delta M \Delta^{-1}\|_p \mid \Delta = \mathrm{diag}(\delta_1, \ldots, \delta_n),\; \delta_i > 0,\; \forall\, i \in \{1, \ldots, n\}\},
\end{equation}
where $\|\cdot\|_p$ is the matrix $p$-norm.

\begin{conjecture}
For any $M \in \mathbb{R}^{n \times n}$ one has
\begin{equation}
\rho_2(M) \leq \rho_3(M) \leq \cdots \leq \rho_\infty(M).
\end{equation}
\end{conjecture}

A weaker version is:
\begin{conjecture}
For any $M \in \mathbb{R}^{n \times n}$ one has
\begin{equation}
\rho_2(M) \leq \rho_k(M), \quad \forall\, k \in \mathbb{N} \setminus \{0\}.
\end{equation}
\end{conjecture}

Another conjecture is:
\begin{conjecture}
For any $n > 6$, and any $p > 2$ ($n, p \in \mathbb{N}$) there exists $M \in \mathbb{R}^{n \times n}$ such that
\begin{equation}
\rho_p(M) < \rho_{p+1}(M).
\end{equation}
\end{conjecture}

\section*{What is already known}

The main previous works are in \cite{coron2008,coron2015}. In particular it is known that for any $n \in \mathbb{N}$ and any $M \in \mathbb{R}^{n \times n}$,
\begin{equation}
\rho_1(M) \geq \rho_2(M), \quad \rho_1(M) = \rho_\infty(M),
\end{equation}
and for $n \geq 2$ there exists $M \in \mathbb{R}^{n \times n}$ such that
\begin{equation}
\rho_1(M) > \rho_2(M).
\end{equation}
It is also known \cite{rump2003} that if $M$ has only non-negative coefficients then
\begin{equation}
\rho_k(M) = \rho_1(M), \quad \forall\, k \in [1, +\infty].
\end{equation}

\section*{Task: Prove Conjecture 3}

We aim to prove:

\medskip
\noindent\textbf{Conjecture 3.} \textit{For any $n > 6$, and any $p > 2$ ($n, p \in \mathbb{N}$) there exists $M \in \mathbb{R}^{n \times n}$ such that $\rho_p(M) < \rho_{p+1}(M)$.}

\newtheorem{lemma}{Lemma}
\newtheorem{theorem}{Theorem}

\section*{Proof of Conjecture 3}

The proof proceeds in three steps. We first compute the induced $\ell_r$ operator norm of the $2\times 2$ Hadamard matrix, then show that diagonal scaling cannot reduce it, and finally observe that zero-padding preserves $\rho_r$. These combine to yield an explicit construction.

\subsection*{Step 1: The induced $\ell_r$ norm of the Hadamard block}

\begin{lemma}\label{lem:hadamard-norm}
Let $H = \begin{pmatrix} 1 & 1 \\ 1 & -1 \end{pmatrix}$. For every $r \in [2,\infty)$,
\[
\|H\|_r = 2^{1-1/r},
\]
where $\|H\|_r = \sup_{x \neq 0} \frac{\|Hx\|_{\ell_r}}{\|x\|_{\ell_r}}$ is the induced $\ell_r$ operator norm.
\end{lemma}

\begin{proof}
Fix $r \in [2,\infty)$. We first record the standard norm comparisons on $\mathbb{R}^2$: for every $z=(z_1,z_2)\in\mathbb{R}^2$,
\[
\|z\|_r \le \|z\|_2
\]
and
\[
\|z\|_2 \le 2^{1/2-1/r}\|z\|_r.
\]
The first inequality follows from monotonicity of finite-dimensional $\ell_p$ norms when $r\ge 2$. For the second, if $r=2$ it is equality, while if $r>2$, H\"older's inequality gives
\[
\|z\|_2^2
=\sum_{j=1}^2 |z_j|^2
\le \left(\sum_{j=1}^2 |z_j|^r\right)^{2/r}
\left(\sum_{j=1}^2 1^{r/(r-2)}\right)^{(r-2)/r}
=2^{1-2/r}\|z\|_r^2.
\]
Taking square roots yields $\|z\|_2 \le 2^{1/2-1/r}\|z\|_r$.

Now let $x=(a,b)\in\mathbb{R}^2$. Then $Hx=(a+b,a-b)$, so
\[
\|Hx\|_2^2 = |a+b|^2+|a-b|^2 = 2a^2+2b^2 = 2\|x\|_2^2.
\]
Hence $\|Hx\|_2=\sqrt{2}\,\|x\|_2$. Using the norm comparisons above, we obtain
\[
\|Hx\|_r \le \|Hx\|_2 = \sqrt{2}\,\|x\|_2 \le \sqrt{2}\,2^{1/2-1/r}\|x\|_r = 2^{1-1/r}\|x\|_r.
\]
Therefore $\|H\|_r\le 2^{1-1/r}$.

It remains to prove sharpness. Take $x=(1,1)$. Then $Hx=(2,0)$, and so
\[
\frac{\|Hx\|_r}{\|x\|_r} = \frac{\|(2,0)\|_r}{\|(1,1)\|_r} = \frac{2}{2^{1/r}} = 2^{1-1/r}.
\]
Thus $\|H\|_r\ge 2^{1-1/r}$. Combining the upper and lower bounds gives $\|H\|_r=2^{1-1/r}$.
\end{proof}

\subsection*{Step 2: Diagonal scaling does not reduce the Hadamard block norm}

\begin{lemma}\label{lem:rho-hadamard}
Let $H = \begin{pmatrix} 1 & 1 \\ 1 & -1 \end{pmatrix}$ and define
\[
\rho_r(H) = \inf\{\|DHD^{-1}\|_r : D = \operatorname{diag}(a,b),\; a,b>0\}.
\]
Then $\rho_r(H) = 2^{1-1/r}$ for every $r \in [2,\infty)$.
\end{lemma}

\begin{proof}
Fix $r \in [2,\infty)$. Let $D=\operatorname{diag}(a,b)$ with $a,b>0$, and set $t=a/b>0$. Then
\[
DHD^{-1} = \begin{pmatrix} 1 & t \\ t^{-1} & -1 \end{pmatrix} =: A_t.
\]
We prove that $\|A_t\|_r \geq 2^{1-1/r}$ for every $t>0$.

We use the elementary fact that for any real numbers $u,v$,
\[
(|u|^r+|v|^r)^{1/r} \geq |u| \quad\text{and}\quad (|u|^r+|v|^r)^{1/r} \geq |v|,
\]
since $|u|^r+|v|^r \geq |u|^r$ and $s \mapsto s^{1/r}$ is increasing.

\textbf{Case $t\geq 1$.} Let $x=2^{-1/r}(1,1)^T$. Then $\|x\|_r = 2^{-1/r}\cdot 2^{1/r} = 1$ and
\[
A_t x = 2^{-1/r}\begin{pmatrix} 1+t \\ t^{-1}-1 \end{pmatrix}.
\]
By the coordinate domination property,
\[
\|A_t x\|_r = 2^{-1/r}(|1+t|^r+|t^{-1}-1|^r)^{1/r} \geq 2^{-1/r}|1+t| \geq 2^{-1/r}\cdot 2 = 2^{1-1/r}.
\]

\textbf{Case $0<t\leq 1$.} Let $y=2^{-1/r}(1,-1)^T$. Then $\|y\|_r = 1$ and
\[
A_t y = 2^{-1/r}\begin{pmatrix} 1-t \\ t^{-1}+1 \end{pmatrix}.
\]
By coordinate domination,
\[
\|A_t y\|_r \geq 2^{-1/r}|t^{-1}+1| \geq 2^{-1/r}\cdot 2 = 2^{1-1/r}.
\]

In both cases, $\|A_t\|_r \geq 2^{1-1/r}$. Taking the infimum over all $t>0$ gives $\rho_r(H)\geq 2^{1-1/r}$.

For the reverse inequality, choose $D=I$:
\[
\rho_r(H) \leq \|H\|_r = 2^{1-1/r},
\]
where the last equality is Lemma~\ref{lem:hadamard-norm}. Combining gives $\rho_r(H)=2^{1-1/r}$.
\end{proof}

\subsection*{Step 3: Zero-block padding preserves $\rho_r$}

\begin{lemma}\label{lem:zero-padding}
Let $B \in \mathbb{R}^{m \times m}$, let $k \in \mathbb{N}$, and let $M = B \oplus 0_k \in \mathbb{R}^{(m+k)\times(m+k)}$. Then $\rho_r(M)=\rho_r(B)$ for every $r \in [1,\infty)$.
\end{lemma}

\begin{proof}
Let $\mathcal{D}_n$ denote the set of positive diagonal matrices in $\mathbb{R}^{n \times n}$. We first establish a basic fact about induced $\ell_r$ operator norms of block diagonal matrices. If $A \in \mathbb{R}^{p \times p}$ and $C \in \mathbb{R}^{q \times q}$, then for every $r \in [1,\infty)$,
\[
\|A \oplus C\|_r=\max\{\|A\|_r,\|C\|_r\}.
\]
Indeed, for vectors $x \in \mathbb{R}^p$ and $y \in \mathbb{R}^q$,
\[
\|(A \oplus C)(x,y)\|_r^r = \|Ax\|_r^r+\|Cy\|_r^r \le \max\{\|A\|_r,\|C\|_r\}^r\bigl(\|x\|_r^r+\|y\|_r^r\bigr),
\]
so $\|A \oplus C\|_r \le \max\{\|A\|_r,\|C\|_r\}$. Conversely, taking $y=0$ gives $\|A \oplus C\|_r \ge \|A\|_r$, and taking $x=0$ gives $\|A \oplus C\|_r \ge \|C\|_r$.

Now let $M=B \oplus 0_k$. Every $D \in \mathcal{D}_{m+k}$ decomposes uniquely as $D=E \oplus F$ where $E \in \mathcal{D}_m$ and $F \in \mathcal{D}_k$. Hence
\[
DMD^{-1} = (E \oplus F)(B \oplus 0_k)(E^{-1}\oplus F^{-1}) = (EBE^{-1}) \oplus 0_k.
\]
Using the block diagonal norm identity,
\[
\|DMD^{-1}\|_r = \max\{\|EBE^{-1}\|_r,\|0_k\|_r\} = \|EBE^{-1}\|_r.
\]
Therefore,
\[
\rho_r(M) = \inf_{D \in \mathcal{D}_{m+k}} \|DMD^{-1}\|_r = \inf_{E \in \mathcal{D}_m} \|EBE^{-1}\|_r = \rho_r(B). \qedhere
\]
\end{proof}


\subsection*{Step 4: Conclusion --- Proof of Conjecture 3}

\begin{theorem}\label{thm:conj3}
For any integers $n > 6$ and $p > 2$, there exists $M \in \mathbb{R}^{n \times n}$ such that $\rho_p(M) < \rho_{p+1}(M)$.
\end{theorem}

\begin{proof}
Let $H = \begin{pmatrix} 1 & 1 \\ 1 & -1 \end{pmatrix}$ and define $M = H \oplus 0_{n-2} \in \mathbb{R}^{n \times n}$, where $0_{n-2}$ denotes the $(n-2)\times(n-2)$ zero matrix. Since $n > 6$, we have $n-2 \geq 5 \geq 0$, so the construction is valid.

By Lemma~\ref{lem:zero-padding} (with $B=H$, $m=2$, $k=n-2$),
\[
\rho_p(M) = \rho_p(H) \quad\text{and}\quad \rho_{p+1}(M) = \rho_{p+1}(H).
\]
By Lemma~\ref{lem:rho-hadamard}, since $p \geq 3 > 2$ and $p+1 \geq 4 > 2$,
\[
\rho_p(H) = 2^{1-1/p} \quad\text{and}\quad \rho_{p+1}(H) = 2^{1-1/(p+1)}.
\]
Since $p < p+1$, we have $\frac{1}{p} > \frac{1}{p+1}$, hence $1 - \frac{1}{p} < 1 - \frac{1}{p+1}$. Because the exponential function $\alpha \mapsto 2^\alpha$ is strictly increasing,
\[
\rho_p(M) = 2^{1-1/p} < 2^{1-1/(p+1)} = \rho_{p+1}(M).
\]
This completes the proof of Conjecture~3.
\end{proof}

\begin{thebibliography}{9}
\bibitem{coron2015} Jean-Michel Coron and Hoai-Minh Nguyen. \textit{Dissipative boundary conditions for nonlinear 1-D hyperbolic systems: sharp conditions through an approach via time-delay systems.} SIAM J.\ Math.\ Anal., 47(3):2220--2240, 2015.

\bibitem{coron2008} Jean-Michel Coron, Georges Bastin, and Brigitte d'Andr\'ea-Novel. \textit{Dissipative boundary conditions for one-dimensional nonlinear hyperbolic systems.} SIAM Journal on Control and Optimization, 47(3):1460--1498, 2008.

\bibitem{rump2003} Siegfried M.\ Rump. \textit{Optimal scaling for $p$-norms and componentwise distance to singularity.} IMA Journal of Numerical Analysis, 23(1):1--9, 2003.
\end{thebibliography}

\end{document}
